\begin{lemma}
For every $\eta>0$ there is $D_0$ such that, whenever $D\ge D_0$, there is a
$3$-uniform $3$-simple hypergraph $\mathcal H$ of maximum degree at most $D$
and an edge $e\in E(\mathcal H)$ with
\[
b_{L(\mathcal H),3D}(e)\ge
1-\frac19+\frac1{3^5}-\eta .
\]
\end{lemma}
\begin{proof}
Choose an integer $D_0$ so large that $D_0\ge 4$ and
$4/D_0\le \eta$.  Fix $D\ge D_0$, and put
\[
q=\left\lfloor\frac{D-1}{3}\right\rfloor .
\]
We construct the required multihypergraph explicitly.  The vertex set is
\[
V=\{0,1,2\}\times\{0,1,2\},
\]
with all arithmetic in the two coordinates taken modulo $3$.  Let
\[
e=\{(0,0),(0,1),(0,2)\}
\]
be the distinguished edge.  For each slope $m\in\{0,1,2\}$ and intercept
$a\in\{0,1,2\}$ define
\[
L_{m,a}=\{(x,mx+a):x\in\{0,1,2\}\}.
\]
The edge-index set consists of one index for $e$ and, for every triple
$(m,a,i)$ with $m,a\in\{0,1,2\}$ and $1\le i\le q$, one further index whose
edge is $L_{m,a}$.  This is a multihypergraph in the sense of
\noderef{MultiHypergraph}: different indices with the same $(m,a)$ and
different $i$ have the same assigned set of vertices.

Each set $L_{m,a}$ has exactly three vertices, because its first coordinates
are the three distinct residues $0,1,2$, and $e$ plainly has three vertices.
Thus every edge has size $3$, so $\mathcal H$ is $3$-uniform by
\noderef{UniformHypergraph}.  Since every edge has size $3$, the intersection
of any two distinct edge indices has size at most $3$; this includes parallel
copies of the same $L_{m,a}$, whose intersection has size exactly $3$.
Therefore $\mathcal H$ is $3$-simple by \noderef{TSimpleHypergraph}.

We next compute degrees, using \noderef{HypergraphDegree}.  A point
$(0,y)$ of $e$ lies in $e$.  For each slope $m$ there is exactly one intercept,
namely $a=y$, for which $(0,y)\in L_{m,a}$, and there are $q$ indices with
that slope and intercept.  Hence
\[
\deg_{\mathcal H}(0,y)=1+3q\le D,
\]
because $3q\le D-1$ by the definition of the floor.  Now let $x\ne 0$.
For each slope $m$ there is exactly one intercept $a=y-mx$ for which
$(x,y)\in L_{m,a}$, and the point is not in $e$.  Thus
\[
\deg_{\mathcal H}(x,y)=3q\le D.
\]
So \noderef{MaxDegreeAtMost} holds with parameter $D$, and together with
the preceding uniformity and simplicity checks, \noderef{HypergraphClass}
places $\mathcal H$ in $H(3,3,D)$.

By \noderef{LineGraphOfHypergraph}, the neighbors of the edge-index $e$ are
exactly the indices whose assigned edge intersects $e$.  For every
$(m,a,i)$, the line $L_{m,a}$ contains the point $(0,a)$ of $e$, so every one
of the $9q$ nonvertical line indices is a neighbor of $e$.  There are no other
edge indices besides $e$ itself.  Hence
\[
\deg_{L(\mathcal H)}(e)=9q.
\]

We record the elementary intersection facts for the sets $L_{m,a}$.  First,
$L_{m,a}\cap e=\{(0,a)\}$, since a point of $e$ has first coordinate $0$, and
the defining equation of $L_{m,a}$ then gives second coordinate $a$.  Second,
if $m$ is fixed and $a\ne b$, then $L_{m,a}$ and $L_{m,b}$ are disjoint:
a common point with first coordinate $x$ would give
$mx+a=mx+b$ modulo $3$, hence $a=b$, a contradiction.  Third, if
$m\ne n$, then $L_{m,a}$ and $L_{n,b}$ meet in exactly one point.  Indeed,
a common point must satisfy $(m-n)x=b-a$ modulo $3$; since the nonzero
residues modulo $3$ are invertible, this equation has exactly one solution
$x\in\{0,1,2\}$, and the corresponding $y=mx+a$ gives the unique common
point.

Using \noderef{IndependentPairCount}, an independent pair in the
neighborhood of $e$ is exactly an unordered pair of neighboring edge-indices
whose assigned hyperedges are disjoint.  The intersection facts show that this
happens exactly when the two indices have the same slope and different
intercepts: equal slopes with different intercepts are disjoint, equal slopes
with the same intercept are parallel copies of the same three-vertex edge and
therefore intersect, and different slopes intersect in one point.  There are
$3$ choices of the slope, $\binom32=3$ unordered choices of the two
intercepts, and then $q$ choices of the copy index over each intercept.
Consequently
\[
P_{L(\mathcal H)}(e)=3\binom32 q^2=9q^2.
\]

Similarly, by \noderef{IndependentTripleCount}, an independent triple in the
neighborhood of $e$ is exactly a three-element set of neighboring edge-indices
that are pairwise disjoint.  The pairwise classification above implies that
all three indices must have the same slope and pairwise different intercepts.
For a fixed slope, this forces one index over each of the three intercepts,
and there are $q^3$ choices of their copy indices.  With $3$ possible slopes,
\[
T_{L(\mathcal H)}(e)=3q^3.
\]

Substituting these exact values into \noderef{LocalBParameter} with scale
$3D$ gives
\[
b_{L(\mathcal H),3D}(e)
=\frac{9q}{3D}-\frac{9q^2}{(3D)^2}+\frac{3q^3}{(3D)^3}
=3\frac qD-\left(\frac qD\right)^2+\frac19\left(\frac qD\right)^3.
\]
Write $x=q/D$ and $c=1/3$.  Since $q=\lfloor(D-1)/3\rfloor$, there is
$s\in\{1,2,3\}$ with $D=3q+s$, so
\[
0\le c-x=\frac{s}{3D}\le \frac1D .
\]
For $0\le x\le c$, the polynomial $f(t)=3t-t^2+t^3/9$ satisfies
\[
f(c)-f(x)=(c-x)\left(3-(c+x)+\frac{c^2+cx+x^2}{9}\right).
\]
The parenthesized factor is less than $4$, because $0\le x,c\le 1$ makes it
at most $3+1/3$.  Hence
\[
\left(1-\frac19+\frac1{3^5}\right)
-b_{L(\mathcal H),3D}(e)
=f(c)-f(x)
< \frac4D
\le \frac4{D_0}
\le \eta .
\]
Therefore
\[
b_{L(\mathcal H),3D}(e)\ge
1-\frac19+\frac1{3^5}-\eta,
\]
as required.
\end{proof}
